%设置页面边距（word标准页面）
\documentclass[a4paper]{article}
\usepackage{geometry}
\geometry{a4paper,left=2.7cm,right=2.7cm,top=2.54cm,bottom=2.54cm}
 
%导入ctex包
\usepackage[UTF8,heading=true]{ctex}

% \usepackage{graphicx}
% \usepackage{caption}
% \usepackage{subcaption}


%需使用包
\usepackage{svg} 

%导入imakeidx包
\usepackage{imakeidx}
\makeindex % 初始化索引处理
 
%设置摘要格式
\usepackage{abstract}
\setlength{\abstitleskip}{0em}
\setlength{\absleftindent}{0pt}
\setlength{\absrightindent}{0pt}
\setlength{\absparsep}{0em}
\renewcommand{\abstractname}{\textbf{\zihao{4}{摘要部分}}}
\renewcommand{\abstracttextfont}{\zihao{-4}} %设置摘要正文字号
 
%设置页眉和页脚，只显示页脚居中页码
\usepackage{fancyhdr}
\pagestyle{plain}
 
%调用数学公式包
\usepackage{amssymb}
\usepackage{amsmath}
 
%调用浮动包
\usepackage{float}
\usepackage{subfig}
\captionsetup[figure]{labelsep=space} %去除图标题的冒号
\captionsetup[table]{labelsep=space} %去除表格标题的冒号
%设置标题格式
\ctexset {
	%设置一级标题的格式
	section = {
		name={,、},
		number=\chinese{section}, %设置中文版的标题
		aftername=,
	},
	%设置三级标题的格式
	subsubsection = {
		format += \zihao{-4} % 设置三级标题的字号
	}
}
 
 
%使得英文字体都为Time NewTown
%\usepackage{times}
 
%图片包导入
% \usepackage{graphicx}
% \graphicspath{{figures}} %图片在当前目录下的figures目录

 
%参考文献包引入
\usepackage{cite}
\usepackage[numbers,sort&compress]{natbib}

 
%代码格式
\usepackage{listings}
\usepackage{graphicx}%写入python代码
\usepackage{pythonhighlight}%python代码高亮显示
\lstset{
	%numbers=left, %设置行号位置
	%	numberstyle=\tiny, %设置行号大小
	keywordstyle=\color{blue}, %设置关键字颜色
	commentstyle=\color[cmyk]{1,0,1,0}, %设置注释颜色
	escapeinside=``, %逃逸字符(1左面的键)，用于显示中文
	breaklines, %自动折行
	extendedchars=false, %解决代码跨页时，章节标题，页眉等汉字不显示的问题
	xleftmargin=1em,xrightmargin=1em, aboveskip=1em, %设置边距
	tabsize=4, %设置tab空格数
	showspaces=false %不显示空格
}
 
 
\renewcommand{\refname}{}
 
%item包
\usepackage{enumitem}
 
%表格加粗
\usepackage{booktabs}
 
%设置表格间距
\usepackage{caption}
 
%允许表格长跨页
\usepackage{longtable}
 
%伪代码用到的宏包
\usepackage{algorithmic}
\usepackage{algorithm}
 
%添加hyperref宏包
\usepackage{hyperref}



%正文区
\title{\Huge\bfseries 2026年信电学院电子设计竞赛}
\date{\LARGE\bfseries 大二六组串口图像检测系统报告} %不显示日期

%文档
\begin{document}
	\maketitle
	\vspace{10em} %设置摘要与标题的间距
	\zihao{-4} %设置正文字号

	\vspace{1cm}
	\vspace{1cm}
	\vspace{1cm}
	\vspace{1cm}
	\begin{table}[ht]
		\centering  
	\label{table:example}  % 为表格提供标签，以便于引用
\end{table}
	\newpage
	\tableofcontents % 插入目录


	\newpage

\section{测试方案与测试结果}


\subsection{模拟检测测试}

测试不同JPEG文件大小情况下LED显示情况。


\begin{table}[H]
\centering
\caption{文件帧长度检测结果}
\begin{tabular}{cccc}
\toprule
测试图片&文件大小(Byte)&理论等级&实际LED显示\\
\midrule
1&&&\\
2&&&\\
3&&&\\
4&&&\\
\bottomrule
\end{tabular}
\end{table}



\subsection{STM32图像显示测试}


\begin{table}[H]
\centering
\caption{图像显示测试结果}
\begin{tabular}{ccc}
\toprule
测试项目&指标&结果\\
\midrule
图像显示&是否正常&\\
帧率&fps&\\
熵值计算&是否准确&\\
\bottomrule
\end{tabular}
\end{table}



\subsection{闭眼检测测试}


\begin{table}[H]
\centering
\caption{闭眼检测结果}
\begin{tabular}{cccc}
\toprule
状态&文件大小变化&LED显示&检测结果\\
\midrule
睁眼&&&\\
闭眼&&&\\
\bottomrule
\end{tabular}
\end{table}



\newpage
\section*{附录}

附录1：模拟检测电路原理图

附录2：PCB及实物照片

附录3：系统整体框图

		

	
\end{document}